\sectionkicker{Theme synthesis}
\subsection*{横断比較 — 2025年後期をExecution Stackとして読み直す}
\addcontentsline{toc}{subsection}{横断比較 — 2025年後期をExecution Stackとして読み直す}
半年を通して一貫して増えたのはmodel名だけではない。model family、agent execution、open deployment、multimodal/action space、runtime、controlがそれぞれ別のchronologyを持ちながら依存し始めたため、技術動向の観測単位をstackへ広げる必要がある。
\medskip
\begin{center}
\small
\renewcommand{\arraystretch}{1.16}
\begin{tabularx}{\linewidth}{@{}p{0.20\linewidth}X@{}}
\toprule
観察軸 & 一次資料・論文から読めること \\
\midrule
7〜8月: layerの立ち上がり & ChatGPT agent、GPT-5、gpt-ossなどはaction surface、frontier family、open-weight deploymentという異なるlayerを早い段階で前景化した。 \cite{src-7a7213b68631,src-b074ccdb2d9e,src-aec54e92a7f3,src-05738f6c55b7,src-964121647555,src-1c46fdf17c79,src-12865e8612c7} \\
9〜10月: 実行系の接続 & agent/coding platform、browser/computer-use、runtime integrationが重なり、model capabilityを実際のactionへ接続するharnessとdownstream implementationの重要性が増した。 \cite{src-f6aa679babd7,src-43eaa2f12960,src-58ef6e84aa65,src-455a78bc36fa} \\
後半: family churnとcontrol & Anthropicのmodel family更新が続く一方、cross-lab safety evaluationなどcontrol layerも独立したengineering concernとして前景化した。 \cite{src-f6aa679babd7,src-ad13fe7d3321} \\
identityを分ける必要 & model announcement、API/service availability、product rollout、runtime integrationは同じ日付・同じ意味ではない。半年を正しく比較するにはこのidentity separationが必須になった。 \cite{src-a24fd97eb3d0,src-43eaa2f12960,src-58ef6e84aa65,src-455a78bc36fa} \\
2025年末に残った境界 & open weightsと全面的なtraining transparency、research prototypeとGA、runtime supportとproduction validation、controlled safety evaluationとdeployed behaviorは依然として別物である。stack化は境界を消すのではなく、むしろ明示する必要を増やした。 \cite{src-1c46fdf17c79,src-12865e8612c7,src-dd5d9b8662f3,src-43eaa2f12960,src-58ef6e84aa65,src-455a78bc36fa,src-ad13fe7d3321} \\
\bottomrule
\end{tabularx}
\renewcommand{\arraystretch}{1.0}
\end{center}
\normalsize
\begin{claimboundary}[この比較の境界]
この表は本号で既に検証・参照している一次資料と論文を横断比較の形へ再配置したものである。vendor / project / author が報告した性能値や優位性は独立再現済みの一般則へ変換せず、本文とSource Notesの帰属境界をそのまま維持する。
\end{claimboundary}
